\chapter{Literature Review}
\lhead{\emph{Chapter 2}}

This chapter surveys the body of knowledge that underpins the proposed framework for
graph-based fraud detection with reinforcement learning and large language model
assistance. The review is organized into six thematic areas: the fraud detection
landscape and its core challenges; classical and machine learning approaches; deep
learning methods; graph-based and graph neural network (GNN) approaches; reinforcement
learning applied to graphs; and the emerging integration of large language models with
graph learning. Each section traces the development of the field, identifies key
limitations that motivated subsequent research, and situates the present work within the
broader literature.

% ─────────────────────────────────────────────────────────────────────────────
\section{The Fraud Detection Landscape}
% ─────────────────────────────────────────────────────────────────────────────

Fraud is a persistent and costly phenomenon that spans virtually every digital and
financial domain. Credit card fraud, healthcare billing abuse, insurance fraud, and
opinion spam on e-commerce platforms collectively generate losses estimated in the
trillions of dollars annually \cite{ngai2011application}.
\cite{bolton2002statistical} provided one of the earliest systematic surveys of
statistical fraud detection, establishing the foundational taxonomy of supervised,
unsupervised, and semi-supervised detection strategies. \cite{chandola2009anomaly}
subsequently broadened this perspective with a comprehensive review of anomaly detection
techniques across domains, noting that fraud detection is distinguished by the adaptive,
adversarial nature of the anomalies—fraudsters continuously modify their behavior in
response to deployed countermeasures.

Three structural characteristics make fraud detection substantially harder than standard
binary classification. First, class imbalance is extreme: fraudulent events typically
represent fewer than one percent of all transactions \cite{dalpozzolo2015calibrating}.
Second, the cost of errors is asymmetric; missing a fraudulent transaction (false
negative) generally incurs far greater cost than a false alarm (false positive)
\cite{baesens2015fraud}. Third, fraudsters exhibit \emph{camouflage} behavior,
deliberately manipulating observable features and social connections to resemble
legitimate entities \cite{dou2020enhancing}. These three challenges motivate the
transition from simple rule-based systems to increasingly sophisticated learning
frameworks.

% ─────────────────────────────────────────────────────────────────────────────
\section{Classical and Machine Learning Approaches}
% ─────────────────────────────────────────────────────────────────────────────

\subsection{Rule-Based and Statistical Methods}

The earliest automated fraud detection systems relied on hand-coded logical rules that
flagged transactions exceeding fixed thresholds or matching known fraud signatures
\cite{phua2010comprehensive}. While operationally simple, such systems are rigid:
they cannot generalize to novel fraud patterns and require continuous manual maintenance
as fraudster behavior evolves. Statistical methods introduced greater analytical rigor.
\cite{bolton2002statistical} applied peer group analysis using $t$-statistics to detect
behavioral deviations in cardholder accounts without requiring labeled fraud examples.
Benford's Law—the prediction that leading digits in naturally occurring numerical data
follow a logarithmic distribution—was adapted as an anomaly signal for insurance and
financial statement fraud \cite{maher2002benford, lu2005detecting}. Although these
approaches can operate without labels, they lack the discriminative power needed for
precision-critical deployment.

\subsection{Supervised Machine Learning}

Supervised learning marked a significant step forward by enabling models to learn
decision boundaries from labeled examples. Logistic regression became the standard
interpretable baseline for binary fraud classification \cite{ngai2011application}.
Decision trees and random forests demonstrated strong performance across systematic
benchmarks; \cite{whitrow2009transaction} found that random forests outperformed
support vector machines, $k$-nearest neighbors, and naïve Bayes on credit card fraud
detection tasks when combined with transaction-aggregation features computed over
rolling time windows. \cite{bhattacharyya2011data} conducted a comparative study
confirming that support vector machines and random forests were the most consistently
competitive classical methods, with SVMs excelling on high-dimensional feature spaces.
Cost-sensitive learning addresses the asymmetric error cost directly: \cite{sahin2013cost}
showed that embedding misclassification costs into a decision tree's splitting criterion
substantially improved operational savings compared to accuracy-optimized variants.

More recently, gradient boosting ensembles such as XGBoost and LightGBM have come
to dominate competition benchmarks and production systems alike. Stacking ensembles
combining multiple diverse base learners have achieved $F_1$-scores exceeding 88\% on
standard credit card fraud benchmarks \cite{ileberi2024enhancing}. Despite these
advances, classical approaches share a fundamental limitation: they operate on fixed,
hand-crafted feature representations and cannot directly incorporate relational
structure among entities—a critical shortcoming given that fraud is inherently a
networked phenomenon.

\subsection{Feature Engineering Strategies}

The quality of input representations profoundly influences detection performance
irrespective of model sophistication.\cite{bahnsen2016feature} introduced periodic features modeled via the
von Mises distribution to capture cyclical spending rhythms, reporting a 13\% average
improvement in cost savings over baseline feature sets. \cite{dalpozzolo2014learned}
demonstrated that incorporating recent consumer behavior improves classifier performance
by over 200\% compared to raw transaction features alone, underscoring the importance
of recency-weighted aggregation. \cite{lucas2019towards} applied multi-perspective
Hidden Markov Models to generate likelihood-based features from transaction sequences,
which improved random forest performance by providing a compact summary of temporal
dynamics. On the network side, \cite{vanvlasselaer2015apate} showed that integrating
personalized PageRank scores computed on tripartite cardholder-transaction-merchant
graphs substantially increased detection performance, providing an early demonstration
that relational features complement tabular ones.

% ─────────────────────────────────────────────────────────────────────────────
\section{Deep Learning Approaches}
% ─────────────────────────────────────────────────────────────────────────────

Deep learning transformed fraud detection by enabling models to learn hierarchical,
nonlinear representations directly from raw data, reducing reliance on manual feature
engineering.

\subsection{Convolutional and Recurrent Architectures}

Convolutional Neural Networks (CNNs) were adapted for fraud detection by representing
transaction feature matrices as pseudo-images amenable to convolutional processing.
\cite{fu2016credit} proposed an early CNN framework that extracts latent patterns from
credit card transaction data, demonstrating performance superior to shallow classifiers
on real-world bank data. Hybrid architectures combining CNNs with SVMs—where the CNN
performs feature extraction and the SVM performs classification—were subsequently shown
to outperform purely convolutional models, suggesting complementarity between
representational and discriminative learning \cite{berhane2023hybrid}.

Recurrent Neural Networks (RNNs), particularly Long Short-Term Memory (LSTM) networks,
addressed the inherently sequential nature of financial transactions. \cite{jurgovsky2018sequence}
framed fraud detection as a sequence classification problem, demonstrating that LSTMs
improved detection of card-present fraud and, critically, detected different fraud types
than random forests, suggesting that ensemble combination of sequential and tabular
models is beneficial. \cite{benchaji2021enhanced} augmented LSTM with attention
mechanisms, showing that attention weights selectively focus on the most fraud-predictive
events within a transaction history while simultaneously providing interpretability for
auditors.

\subsection{Autoencoders and Generative Models}

Autoencoders offered an unsupervised route to anomaly detection that is well suited to
the extreme class imbalance in fraud settings. Trained exclusively on legitimate
transactions, an autoencoder learns to reconstruct normal patterns and assigns high
reconstruction error to anomalous inputs. \cite{schreyer2017detection} applied deep
autoencoders to detect anomalous journal entries in financial audits, and their
follow-up work \cite{schreyer2019adversarial} extended this to adversarial autoencoders
that learn semantically interpretable latent spaces, providing both detection and
explanation capabilities.

Generative Adversarial Networks (GANs) addressed class imbalance by generating
synthetic fraudulent examples to augment training sets. \cite{fiore2019using}
demonstrated that classifiers trained on GAN-augmented datasets significantly improved
sensitivity, and \cite{alamri2023survey} surveyed GAN variants for this purpose,
finding that Wasserstein GANs and Conditional Tabular GANs are particularly effective
for the mixed-type tabular data common in fraud datasets.

\subsection{Transformer and Attention-Based Methods}

Transformer architectures represent the most recent deep learning frontier in fraud
detection. \cite{yu2024credit} adapted the full Transformer encoder for credit card
fraud detection, reporting significant improvements over XGBoost and TabNet baselines.
The Finsformer model \cite{finsformer2024} integrated Transformers with cluster-based
attention mechanisms, achieving high precision and recall on financial attack detection.
A recurring finding across attention-based studies is that attention weights serve as
natural, post-hoc explanations of model decisions—an important property in regulated
financial environments where model transparency is a compliance requirement.

Despite their representational power, all of these architectures treat each transaction
or entity as an independent sample or as part of a sequence. They cannot directly encode
the relational dependencies between entities that are central to understanding networked
fraud. This limitation motivates the adoption of graph-based representations.

% ─────────────────────────────────────────────────────────────────────────────
\section{Graph-Based Methods and Graph Neural Networks}
% ─────────────────────────────────────────────────────────────────────────────

\subsection{Early Graph-Based Fraud Detection}

The recognition that fraud is fundamentally a relational phenomenon motivated an early
shift toward graph-based analysis. \cite{pandit2007netprobe} introduced NetProbe, which
models auction users and transactions as a Markov Random Field and applies belief
propagation to infer fraud probabilities, establishing the \emph{guilt-by-association}
principle—that suspicion propagates through network connections. \cite{akoglu2015graph}
provided a comprehensive survey of graph-based anomaly detection, covering spectral
methods, dense subgraph mining, and community detection, and demonstrated the breadth of
fraud-relevant structural signals available in relational data. Shallow embedding
methods such as DeepWalk \cite{perozzi2014deepwalk} and node2vec \cite{grover2016node2vec}
generated node representations by simulating random walks on the graph, capturing both
local and global structural proximity. Although effective at encoding topology, these
methods cannot jointly leverage node attributes and graph structure in an end-to-end
fashion, limiting their utility in settings where both feature space and relational
structure are informative.

\subsection{Foundational Graph Neural Network Architectures}

Graph Neural Networks (GNNs) addressed the limitations of shallow methods by enabling
end-to-end learning of node representations that jointly incorporate node features and
neighborhood topology through iterative message passing. \cite{bronstein2017geometric}
framed this development within the broader program of geometric deep learning, arguing
that graphs provide the natural inductive bias for relational data. \cite{gilmer2017neural}
proposed the Message Passing Neural Network (MPNN) framework as a unifying abstraction,
showing that most GNN variants can be expressed as instances of a general
neighborhood-aggregation update rule. \cite{wu2021graph} provided a comprehensive
survey of GNN architectures and applications, tracing the field from spectral to spatial
methods and documenting the empirical conditions under which each class performs best.

The Graph Convolutional Network (GCN) of \cite{kipf2017semi} introduced scalable
spectral-inspired convolutions through a first-order approximation of graph Laplacian
filters, with a propagation rule that scales linearly in the number of edges and enables
semi-supervised learning from a small fraction of labeled nodes. GraphSAGE
\cite{hamilton2017inductive} addressed the transductive limitation of GCN by introducing
inductive representation learning through fixed-size neighborhood sampling and flexible
aggregation operators—mean, max-pooling, and LSTM—enabling generalization to previously
unseen nodes at inference time. The Graph Attention Network (GAT) \cite{velickovic2018graph}
extended this framework with masked self-attention over neighborhoods, learning
differential importance weights for different neighbors and providing a natural mechanism
for distinguishing informative from noisy connections. The Relational GCN (RGCN)
\cite{schlichtkrull2018modeling} extended GCN to multi-relational data by learning
relation-specific weight matrices, directly addressing the heterogeneous edge types
present in real fraud networks. Together, these four architectures form the backbone of
virtually all subsequent GNN-based fraud detection systems.

\subsection{GNN Architectures for Fraud Detection}

The application of GNNs to fraud detection proceeded along several parallel lines.
\cite{liu2018heterogeneous} introduced a heterogeneous GNN that aggregates information
across device, activity, and account subgraphs for malicious account detection on
Alipay, establishing an early and influential multi-graph design pattern. GeniePath
\cite{liu2019geniepath} proposed adaptive receptive paths via LSTM-gated attention,
allowing the model to select the most relevant depth of neighborhood aggregation for each
node. \cite{wang2019semignn} introduced a semi-supervised graph attentive network that
applies hierarchical attention across multiple graph views, improving performance in
settings with limited labels. \cite{li2019spam} designed a heterogeneous GCN for spam
review detection on a commodity marketplace that applies type-specific aggregators for
different relation types.

\cite{liu2020alleviating} formally identified the \emph{inconsistency problem} in
applying GNNs to fraud detection: because fraudsters strategically connect to benign
nodes, standard message passing aggregates predominantly legitimate neighbor signals,
which suppresses the fraudulent signal of the target node. Their GraphConsis model
addresses this by filtering neighbors whose feature vectors are inconsistent with the
center node's class prediction before aggregation. \cite{dou2020enhancing} extended
this line of work by formalizing two distinct types of fraudster camouflage—feature
camouflage, in which fraudsters manipulate their own attributes to resemble benign users,
and relation camouflage, in which they establish connections to legitimate entities—and
proposing a reinforcement learning–guided neighbor selector that adaptively filters
noisy connections at the relation level. \cite{liu2021pick} addressed the class
imbalance dimension of the problem through label-balanced neighborhood sampling, showing
that ensuring class balance during aggregation consistently improves minority-class
recall. \cite{peng2022reinforced} further developed the RL-based approach with a
recursive and scalable neighbor selection framework that operates across multiple
relation types, demonstrating consistent improvements on standard benchmarks.

\subsection{Scalability, Explainability, and Robustness}

Three additional concerns have shaped recent GNN fraud detection research. First,
\emph{scalability}: production fraud systems must operate on graphs with hundreds of
millions of nodes and billions of edges. \cite{wu2019simplifying} showed that removing
nonlinearities between GCN layers and collapsing them into a single linear transformation
preserves most of the expressive power while dramatically reducing computational cost.
xFraud \cite{rao2022xfraud}, deployed at eBay on graphs with over one billion nodes,
demonstrated that heterogeneous GNNs can be scaled to production through efficient
mini-batch sampling and embedding caching. Second, \emph{explainability}: the need for
regulatory transparency in financial services has motivated the development of
interpretable fraud detectors. SEFraud \cite{sefraud2024} generates self-explanations
through interpretive mask learning on graph structures, identifying the subgraph
substructures most responsible for a fraud prediction. Third, \emph{adversarial
robustness}: \cite{sun2018adversarial} surveyed attacks on graph data—including node
feature perturbation and edge injection—and corresponding defenses, documenting that
GNNs are substantially more vulnerable to adversarial manipulation than their Euclidean
counterparts, which is a critical concern in adversarial fraud settings.

% ─────────────────────────────────────────────────────────────────────────────
\section{Reinforcement Learning for Graph-Based Decision Making}
% ─────────────────────────────────────────────────────────────────────────────

Reinforcement learning (RL) provides a principled framework for sequential decision
making under uncertainty and has seen growing application to graph-structured problems.
\cite{zhou2020graph} surveyed GNN methods and applications, highlighting the potential
of combining learned graph representations with RL policies for downstream tasks
including graph construction, node selection, and dynamic network control. The
combination is natural: GNNs provide rich structural representations of system state,
while RL policies learn action strategies that maximize long-term reward, enabling
adaptive behavior that static models cannot achieve.

In the fraud detection context, the key decision problem is \emph{relation weighting}:
which relational channels and which specific neighbors should be trusted when aggregating
messages for a given node? Static approaches assign fixed weights to all relations, which
is suboptimal when the informativeness of each relation varies across nodes and evolves
as fraudsters adapt their camouflage strategies. Early work by \cite{dou2020enhancing}
formulated this as a multi-armed bandit problem, using cumulative reward based on
inter-epoch changes in average neighbor distance to adaptively adjust the top-$p$
filtering threshold for each relation. \cite{peng2022reinforced} refined this with a
recursive RL framework that operates at multiple relation levels simultaneously,
improving both the granularity and stability of the learned selection policy.

More broadly, RL has been applied to graph-based tasks including combinatorial
optimization, knowledge graph completion, and dynamic graph evolution
\cite{zhou2020graphrl}. Policy gradient methods such as REINFORCE, proximal policy
optimization (PPO), and deep $Q$-networks (DQN) have all been applied in graph settings,
with the specific choice depending on whether the action space is continuous or discrete
and whether a differentiable reward signal is available. A general challenge in applying
RL to graph problems is the \emph{state representation}: the RL agent must receive a
compact and informative summary of the current graph state to make good decisions. Early
approaches used handcrafted structural statistics; recent work increasingly uses GNN
embeddings as state representations, closing the loop between the representation learner
and the policy \cite{zhou2020graphrl}.

% ─────────────────────────────────────────────────────────────────────────────
\section{Large Language Models for Semantic Graph Analysis}
% ─────────────────────────────────────────────────────────────────────────────

Large language models (LLMs) have emerged as a powerful tool for extracting semantic
representations from textual data, with recent work beginning to integrate them with
graph neural networks for tasks where both structural and linguistic signals are relevant.
In fraud detection, text associated with entities—review content, transaction
descriptions, user profiles—carries semantic anomaly signals that complement topological
ones. Fraudsters who mimic legitimate relational structure may still exhibit detectable
linguistic patterns, and conversely, textual analysis alone misses the relational context
that reveals coordinated fraud rings.

\cite{yang2025flag} introduced FLAG, the first framework to integrate LLMs with GNNs
for fraud detection. FLAG uses an LLM to extract discriminative text features from
entity descriptions and applies them at two levels: to supplement node feature
representations for the GNN and to guide neighbor sampling by computing LLM-based
semantic similarity between candidate neighbors and the target node. Deployed in Alipay's
credit risk system, FLAG reported average improvements of approximately 3\% in
$F_1$-macro and 7\% in AUC over competitive GNN baselines on public benchmarks,
demonstrating that LLM-derived semantic features provide complementary signal beyond
what is captured by structural aggregation alone.

This integration is conceptually important because it bridges two distinct modes of
reasoning: structural reasoning, which operates over the topology of interaction graphs,
and semantic reasoning, which operates over the meaning of entity descriptions and
behavioral patterns. Prior graph fraud detection systems were limited to the former.
By incorporating LLM-based semantic state representations into the RL policy, as
proposed in this thesis, the agent can detect feature-level camouflage that is invisible
to purely structural analysis—for instance, identifying that a node's textual description
is semantically anomalous relative to its claimed neighborhood, even when its graph
position has been carefully constructed to appear legitimate.

% ─────────────────────────────────────────────────────────────────────────────
\section{Summary and Research Gap}
% ─────────────────────────────────────────────────────────────────────────────

The literature reviewed in this chapter reveals a clear and consistent progression.
Classical and machine learning approaches established effective baselines but were
limited by hand-crafted features and their inability to model relational structure.
Deep learning methods—CNNs, RNNs, autoencoders, transformers—improved representational
power but still treat entities as independent instances or linear sequences, missing
the graph topology. Graph neural networks addressed this gap by enabling end-to-end
learning over relational data, and purpose-built fraud detection architectures added
mechanisms for class imbalance, camouflage resistance, and heterogeneous relations.
Reinforcement learning extended GNN-based detectors by enabling adaptive, policy-driven
decisions about which relational signals to trust. Most recently, large language models
have opened a new dimension by introducing semantic reasoning into what had previously
been a purely structural analysis.

However, a gap remains. Existing RL-guided GNN frameworks either use simple bandit-based
controllers with limited state representations \cite{dou2020enhancing}, or do not
incorporate semantic signals into the policy state at all \cite{peng2022reinforced}.
Conversely, LLM-integration work such as FLAG uses language models to augment node
features and guide sampling, but does not employ a learned RL policy for dynamic
relation weighting. No existing work unifies all three components—GNN-based structural
reasoning, RL-guided relation weighting, and LLM-enhanced semantic state modeling—into
a single coherent framework. This thesis addresses that gap.

\clearpage